\section{从基准饱和到多模态 Agent：问题空间为什么被重写}

\subsection{主讲脉络与核心判断}
本讲由 Salesforce AI Research 的 Caiming Xiong 主讲，主线非常清晰：传统文本任务基准在过去几年接近饱和，下一阶段的挑战不再是 ``离线答题''，而是 ``在真实数字环境中持续完成任务''。这一定义把研究对象从静态 NLP 任务，切换为可交互、多步骤、跨应用的 Agent 系统。

\begin{importantbox}{从 ``会回答'' 到 ``会做事''}
在多模态 Agent 范式里，模型价值不只是输出一句话是否正确，而是能否稳定地完成一个真实工作流：读界面、理解状态、规划动作、执行操作、校验结果、再迭代。这个闭环才是生产环境里的真正门槛。
\end{importantbox}

\subsection{为什么 ``多模态'' 不再是可选项}
讲者把多模态 Agent 定义为 Vision-Language-Action（VLA）系统。原因很直接：现实工作流并不是纯文本 API，更多是 GUI、网页、桌面软件、移动端界面，甚至外设交互。Agent 只有同时具备视觉理解、语言推理和动作生成能力，才可能接管复杂流程。

\begin{knowledgebox}{多模态 Agent 的输入输出对比}
\begin{itemize}
    \item 输入从 ``单轮文本'' 变成 ``用户指令 + 屏幕截图 + 可访问性树 + 环境状态''；
    \item 输出从 ``自然语言答案'' 变成 ``可执行动作序列''；
    \item 评估从 ``准确率'' 变成 ``任务完成率、轨迹质量、步骤效率与鲁棒性''。
\end{itemize}
\end{knowledgebox}

\begin{figure}[H]
\centering
\includegraphics[width=0.88\textwidth]{frame-01-benchmark.jpg}
\caption{视频约 00:01:30：讲者展示近年基准快速饱和，并引出 Agent 评测要从离线题库转向真实环境}
\end{figure}

\begin{warningbox}{常见误解：把 ``会用工具'' 当成 ``可部署 Agent''}
演示里一次成功并不代表系统可上线。生产级 Agent 关注的是长期稳定性：错误恢复、任务中断处理、跨应用状态一致性、权限边界和审计可追踪性。这些都不是单个 benchmark 分数能覆盖的。
\end{warningbox}

\subsection{本章小结}
多模态 Agent 的研究重点已从模型 ``答题能力'' 转向 ``环境交互能力''。这一转变决定了后续必须同时建设三层基础：真实环境 benchmark、可扩展数据管线、可执行的模型架构。

